%! Tex root = main.tex
\documentclass{article}
\usepackage[utf8]{inputenc}
\usepackage[T1]{fontenc}
\usepackage[brazil]{babel}
\usepackage[a4paper, total={6in, 8in}]{geometry}
\usepackage[portuguese, ruled, linesnumbered]{algorithm2e}
\usepackage{amsmath, amssymb}
\usepackage{booktabs}
\usepackage{graphicx}
\usepackage{hyperref}

\title{SPU2D Bin Conflicts:\\ Empacotamento Bidimensional em Faixas com Restrições de Conflitos}
\author{Ycaro Sales}

\begin{document}
\maketitle

\begin{abstract}
Este trabalho aborda o Problema de Empacotamento Bidimensional em Faixas com
Restrições de Conflitos (\textit{2D Strip/Bin Packing with Conflicts}), uma
variante NP-difícil do empacotamento clássico em que pares de itens
incompatíveis não podem ser alocados em um mesmo recipiente. O objetivo é
empacotar todos os itens retangulares no menor número possível de faixas de
dimensões fixas, respeitando a ausência de sobreposição e o grafo de conflitos.
Apresentamos uma formulação do problema, uma representação da solução baseada em
chaves aleatórias e duas estratégias de resolução: uma heurística construtiva
gulosa do tipo \textit{First-Fit-Decreasing} e uma metaheurística
\textit{Biased Random-Key Genetic Algorithm} (BRKGA). Os experimentos sobre 18
instâncias mostram que o BRKGA reduz o número total de faixas de 91 para 78
(redução média de aproximadamente 14{,}3\%), superando a heurística gulosa em 11
instâncias e empatando nas demais, sem nunca produzir resultado pior.
\end{abstract}

\newpage
\tableofcontents
\newpage

\section{Introdução}

A logística ocupa papel central na economia contemporânea: o transporte e o
armazenamento de mercadorias representam parcela expressiva dos custos
operacionais de empresas de manufatura, varejo e comércio eletrônico. Nesse
contexto, problemas de \emph{corte e empacotamento} (\textit{cutting and
packing}) surgem sempre que é preciso acomodar um conjunto de objetos dentro de
recipientes maiores --- caixas em paletes, itens em contêineres, peças em chapas
de matéria-prima --- de forma a aproveitar ao máximo o espaço disponível. Cada
recipiente economizado significa menos viagens, menos área de estoque e menor
desperdício de material, com impacto direto sobre custos e sustentabilidade.

Este trabalho trata de uma generalização do empacotamento bidimensional em que,
além das restrições geométricas usuais, existe a noção de \emph{conflito} entre
itens. Dois itens em conflito não podem compartilhar o mesmo recipiente. Esse
tipo de restrição modela situações reais como produtos químicos incompatíveis,
mercadorias destinadas a clientes ou rotas distintas, ou itens que, por razões
de segurança ou regulação, não podem ser transportados juntos. O acréscimo dos
conflitos torna o problema substancialmente mais difícil, pois reduz a
liberdade de combinação dos itens e tende a aumentar o número de recipientes
necessários.

O objetivo deste trabalho é minimizar o número de faixas (recipientes) utilizadas
para empacotar todos os itens de uma instância. Para isso, implementou-se em
Rust um \textit{solver} que oferece duas abordagens: uma heurística construtiva
gulosa, usada como linha de base, e uma metaheurística evolutiva (BRKGA),
descritas na Seção~\ref{sec:metodologia}. As demais seções estão organizadas da
seguinte forma: a Seção~\ref{sec:problema} caracteriza formalmente o problema; a
Seção~\ref{sec:fundamentacao} apresenta a fundamentação teórica; a
Seção~\ref{sec:metodologia} descreve a metodologia, incluindo a representação da
solução e as heurísticas empregadas; e a Seção~\ref{sec:resultados} reporta os
resultados experimentais.

\section{Caracterização do problema}
\label{sec:problema}

\textbf{Entrada:} um conjunto de faixas (recipientes) retangulares de largura
fixa $W$ e altura fixa $H$; um conjunto finito
$\Re = \{r_{1}, r_{2}, \ldots, r_{n}\}$ de $n$ retângulos, representando os
itens; e um grafo simples de conflitos $G=(\Re, E)$, com
$E\subseteq \{\{u,v\} \mid u,v\in \Re \wedge u \neq v\}$. Cada item $r_{i}$
possui largura $w_{i}$, altura $h_{i}$ e classe de cliente $c_{i}$.

\textbf{Objetivo:} empacotar todos os itens $r\in\Re$ dentro das faixas,
minimizando o número de faixas utilizadas.

\textbf{Restrições:}
\begin{itemize}
    \item Todos os itens devem ser empacotados;
    \item não são permitidas sobreposições entre itens em uma mesma faixa;
    \item cada item deve estar inteiramente contido nos limites da faixa que o
          contém ($0 \le x$, $0 \le y$, $x + w_i \le W$, $y + h_i \le H$),
          admitindo-se rotação de $90^\circ$;
    \item a restrição de conflitos deve ser satisfeita: para toda aresta
          $\{u,v\} \in E$, os itens $u$ e $v$ não podem ser alocados em uma
          mesma faixa.
\end{itemize}

\subsection{O problema do empacotamento bidimensional em faixa com restrições de descarregamento e de conflitos}

O problema estudado é uma variante do empacotamento bidimensional em faixas
(\textit{2D Strip/Bin Packing}). Na formulação clássica, busca-se posicionar
retângulos sem sobreposição dentro de recipientes de dimensão fixa, minimizando
o número de recipientes. Aqui, a essa estrutura acrescentam-se: (i) a
\emph{restrição de conflitos}, modelada pelo grafo $G$, segundo a qual itens
ligados por uma aresta não podem coexistir no mesmo recipiente; e (ii) a
informação de \emph{classe de cliente} $c_i$ de cada item, que dá suporte a uma
\emph{restrição de descarregamento} --- a ordem em que os itens de diferentes
clientes são acomodados pode condicionar a viabilidade prática do descarregamento
ao longo de uma rota. Neste trabalho, o foco recai sobre a restrição de
conflitos e sobre a minimização do número de faixas; a classe de cliente é
representada nos dados de entrada e fica disponível para extensões futuras.

\section{Fundamentação teórica}
\label{sec:fundamentacao}

O empacotamento bidimensional pertence à família de problemas de corte e
empacotamento e é NP-difícil, generalizando o \textit{Bin Packing}
unidimensional. A introdução do grafo de conflitos aproxima o problema de
formulações de coloração de grafos: se a geometria fosse ignorada, alocar itens
a recipientes respeitando os conflitos equivaleria a colorir $G$ com o menor
número de cores, em que cada cor representa um recipiente. Assim, o número
cromático de $G$ é um limitante inferior para o número de faixas necessárias,
pois itens mutuamente em conflito (uma clique em $G$) exigem recipientes
distintos.

Como o problema é intratável na prática para instâncias de tamanho realista,
recorre-se a métodos aproximados. As \emph{heurísticas construtivas} produzem
uma solução em uma única passagem, tomando decisões locais (por exemplo, regras
de ajuste como \textit{First-Fit} e \textit{Best-Fit}); são rápidas, mas podem
ficar presas em soluções de qualidade limitada. As \emph{metaheurísticas}, por
sua vez, exploram o espaço de soluções de forma iterativa e estocástica,
buscando escapar de ótimos locais. Entre elas, os algoritmos genéticos evoluem
uma população de soluções por meio de seleção, recombinação (\textit{crossover})
e mutação.

O \textbf{BRKGA} (\textit{Biased Random-Key Genetic Algorithm}) é uma
metaheurística genética em que cada solução é codificada por um vetor de chaves
aleatórias reais no intervalo $[0,1]$. Um \emph{decodificador} determinístico,
específico do problema, traduz esse vetor em uma solução concreta e calcula seu
custo. Essa separação entre representação (genérica) e decodificação
(específica) permite reutilizar todo o maquinário evolutivo --- seleção,
\textit{crossover} e mutação operam apenas sobre as chaves --- enquanto o
conhecimento do domínio fica concentrado no decodificador.

Para o posicionamento geométrico dos itens dentro de cada recipiente,
utiliza-se a técnica de \emph{Espaços Máximos Vazios} (\textit{Empty Maximal
Spaces}, EMS): o espaço livre de um recipiente é representado por um conjunto de
retângulos vazios maximais; ao inserir um item, os espaços que ele intercepta
são recortados em novos retângulos e os espaços contidos em outros são
descartados, mantendo apenas os maximais.

\section{Metodologia}
\label{sec:metodologia}

\subsection{Visão geral}

O \textit{solver} foi implementado em Rust. Cada instância é lida de um arquivo
JSON (formatos \texttt{Objects/Items} e \texttt{Container/ItemTypes}),
construindo-se a estrutura \texttt{Instance} com os itens, as dimensões da faixa
e o grafo de conflitos indexado pela posição de cada item. Para cada instância
são executadas duas estratégias --- a heurística gulosa e o BRKGA --- e as
soluções resultantes são serializadas em JSON, registrando o número de faixas
utilizadas, os itens não alocados, o tempo de execução e a posição de cada item.

\subsection{Representação da solução}

A solução de uma instância é o conjunto de faixas, cada uma contendo as
posições (origem, dimensões e orientação) dos itens nela alocados. O componente
geométrico central é o \texttt{Container}, que mantém o conjunto de Espaços
Máximos Vazios (EMS) e as colocações já realizadas. A inserção segue a regra
\textit{bottom-left}: dentre os espaços livres que comportam o item, escolhe-se
o de menor origem (primeiro em $y$, depois em $x$); a checagem de conflito
rejeita o item caso algum vizinho seu em $G$ já esteja na faixa.

No BRKGA, a solução é codificada por um vetor de $2n$ chaves aleatórias em
$[0,1]$, interpretado da seguinte forma:
\begin{itemize}
    \item as primeiras $n$ chaves definem a \textbf{ordem de inserção} dos
          itens: os itens são ordenados de forma crescente pela sua chave;
    \item as últimas $n$ chaves definem a \textbf{orientação} de cada item: uma
          chave maior ou igual a um limiar ($0{,}5$) indica que o item é
          rotacionado em $90^\circ$; caso contrário, mantém a orientação
          original.
\end{itemize}
O decodificador percorre os itens na ordem dada e os insere por
\textit{First-Fit} com orientação fixada, abrindo uma nova faixa quando nenhuma
das abertas aceita o item; um item que não cabe nem em uma faixa vazia na sua
orientação forçada é marcado como não alocado.

\subsection{Heurística gulosa (\textit{First-Fit-Decreasing})}

A heurística gulosa serve de linha de base. Os itens são ordenados por área
decrescente e inseridos, um a um, por \textit{First-Fit} nas faixas abertas;
quando nenhuma faixa aberta aceita o item, abre-se uma nova faixa. A inserção é
\emph{ciente de conflitos} e tenta automaticamente as duas orientações do item.
O Algoritmo~\ref{alg:ffd} resume o procedimento.

\begin{algorithm}[h]
\caption{Heurística gulosa First-Fit-Decreasing}
\label{alg:ffd}
\KwIn{instância com itens $\Re$, faixa $W \times H$, grafo $G$}
\KwOut{conjunto de faixas $B$}
$B \leftarrow \emptyset$\;
ordene $\Re$ por área decrescente\;
\ForEach{item $r \in \Re$}{
    $alocado \leftarrow falso$\;
    \ForEach{faixa $b \in B$}{
        \If{$b$ aceita $r$ (sem conflito e com espaço)}{
            insira $r$ em $b$ (regra bottom-left, auto-rotação)\;
            $alocado \leftarrow verdadeiro$; \textbf{break}\;
        }
    }
    \If{$\neg alocado$}{
        abra nova faixa $b'$ e insira $r$ em $b'$\;
        $B \leftarrow B \cup \{b'\}$\;
    }
}
\Return $B$\;
\end{algorithm}

\subsection{Metaheurística BRKGA}

A metaheurística evolui uma população de vetores de chaves aleatórias. A
\emph{função de aptidão} (\textit{fitness}) de um cromossomo é obtida
decodificando-o e calculando
\[
f = \texttt{faixas\_usadas} + \texttt{itens\_nao\_alocados}\cdot (n+1),
\]
a ser \emph{minimizada}. A penalidade $(n+1)$ por item não alocado garante que
qualquer solução inviável seja sempre pior do que qualquer solução viável (que
usa, no máximo, $n$ faixas), de modo que a busca prioriza a viabilidade. O
processo evolutivo, ilustrado no Algoritmo~\ref{alg:brkga}, foi configurado com
população de $100$ indivíduos, critério de parada de $200$ gerações sem melhora,
seleção por torneio, \textit{crossover} uniforme e mutação de gene único.

\begin{algorithm}[h]
\caption{BRKGA para o empacotamento com conflitos}
\label{alg:brkga}
\KwIn{instância, tamanho de população $P$, máximo de gerações estagnadas $S$}
\KwOut{melhor empacotamento encontrado}
inicialize população de $P$ vetores de $2n$ chaves em $[0,1]$\;
\Repeat{$S$ gerações sem melhora}{
    \ForEach{cromossomo}{
        decodifique (ordem $+$ orientação) e empacote por First-Fit\;
        avalie $f = \texttt{faixas} + \texttt{nao\_alocados}\cdot(n+1)$\;
    }
    selecione por torneio; aplique crossover uniforme e mutação\;
    forme a nova geração\;
}
\Return decodificação do melhor cromossomo\;
\end{algorithm}

\section{Resultados}
\label{sec:resultados}

As duas estratégias foram avaliadas sobre 18 instâncias derivadas de conjuntos
de teste de empacotamento bidimensional, às quais foram adicionados conflitos
com diferentes densidades (sufixos \texttt{c0.1}, \texttt{c0.3} e \texttt{c0.5},
indicando a probabilidade de conflito entre pares de itens). A
Tabela~\ref{tab:resultados} compara o número de faixas utilizadas pela
heurística gulosa e pelo BRKGA.

\begin{table}[h]
\centering
\caption{Número de faixas utilizadas: gulosa vs. BRKGA.}
\label{tab:resultados}
\begin{tabular}{lrrr}
\toprule
Instância & $n$ & Gulosa & BRKGA \\
\midrule
BPPSP-n100m18b2d2\_Hard\_c0.1      & 18 & 3 & 3 \\
BPPSP-n100m18b2d2\_Hard\_c0.5      & 18 & 6 & 5 \\
CJCM-E00N23\_HardInfeasible\_c0.3  & 23 & 6 & 5 \\
CJCM-E00N23\_HardInfeasible\_c0.5  & 23 & 8 & 6 \\
CJCM-E02F20\_ManyCalls\_c0.3       & 20 & 5 & 4 \\
CJCM-E02N20\_Hard\_c0.3            & 20 & 6 & 4 \\
CJCM-E02N20\_Hard\_c0.5            & 20 & 7 & 6 \\
CJCM-E03X18\_HardFeasible\_c0.1    & 18 & 3 & 2 \\
CJCM-E03X18\_HardFeasible\_c0.5    & 18 & 7 & 6 \\
\midrule
\textbf{Total (18 instâncias)}     &    & \textbf{91} & \textbf{78} \\
\bottomrule
\end{tabular}
\end{table}

Considerando o conjunto completo de 18 instâncias, a heurística gulosa utilizou
91 faixas no total, enquanto o BRKGA utilizou 78 --- uma redução média de
aproximadamente $14{,}3\%$. O BRKGA produziu soluções estritamente melhores em
11 instâncias, empatou em 7 e em nenhum caso foi pior do que a gulosa. Quanto ao
tempo de execução, a heurística gulosa é praticamente instantânea
($<0{,}01$\,s por instância), ao passo que o BRKGA, por seu caráter iterativo,
gasta tipicamente entre $0{,}07$\,s e $0{,}13$\,s por instância --- um custo
ainda baixo, compatível com seu ganho de qualidade.

Observa-se também que o aumento da densidade de conflitos (de \texttt{c0.1} para
\texttt{c0.5}) eleva consistentemente o número de faixas necessárias, o que é
esperado: quanto mais arestas no grafo de conflitos, menos itens podem
compartilhar um mesmo recipiente. Nesses cenários mais restritos é justamente
onde o BRKGA mais se distancia da heurística gulosa, evidenciando o valor da
busca metaheurística quando as restrições de conflito se intensificam.

\section{Conclusão}

Este trabalho apresentou uma formulação e duas estratégias de resolução para o
empacotamento bidimensional em faixas com restrições de conflitos. A
representação por chaves aleatórias do BRKGA, combinada a um decodificador
\textit{First-Fit} ciente de conflitos e ao gerenciamento de espaços livres via
EMS, mostrou-se eficaz: o BRKGA reduziu o número total de faixas em relação à
heurística gulosa sem nunca piorar nenhuma instância. Como trabalhos futuros,
destacam-se a incorporação efetiva da restrição de descarregamento por classe de
cliente, a avaliação sobre um conjunto maior de instâncias e o ajuste sistemático
dos parâmetros da metaheurística.

\section*{Referências Bibliográficas}

\begin{thebibliography}{9}
\bibitem{goncalves2011} Gonçalves, J. F.; Resende, M. G. C. (2011).
\textit{Biased random-key genetic algorithms for combinatorial optimization}.
Journal of Heuristics, 17(5), 487--525.

\bibitem{lodi2002} Lodi, A.; Martello, S.; Monaci, M. (2002).
\textit{Two-dimensional packing problems: A survey}. European Journal of
Operational Research, 141(2), 241--252.

\bibitem{wascher2007} Wäscher, G.; Haußner, H.; Schumann, H. (2007).
\textit{An improved typology of cutting and packing problems}. European Journal
of Operational Research, 183(3), 1109--1130.
\end{thebibliography}

\end{document}
